\chapter{Appendix: Number Theory Reference}

This appendix gathers, in one place, the number-theoretic machinery that the rest of the book
leans on. It is meant to be read straight through once and then kept open beside the cryptography
chapters as a lookup table. Where the main text proves a result slowly and motivates it with an
Alice/Bob/Eve story, here we state it crisply, prove it when the proof is short and illuminating,
and --- crucially --- work a tiny numeric example so the formula stops being a string of symbols.
Every fact below shows up somewhere in a real protocol: the extended Euclidean algorithm computes
RSA private keys, Euler's theorem proves RSA decryption undoes encryption, the Chinese Remainder
Theorem quadruples RSA signing speed, Miller-Rabin is how your TLS library finds the primes in the
first place, and quadratic residues underlie everything from Rabin encryption to point compression
on elliptic curves.

Nothing here is deep in the sense of requiring a research background. It is deep in the sense that
the same dozen facts, combined in different ways, generate almost all of classical public-key
cryptography. Learn them once, trace them on small numbers, and the schemes stop looking like magic.

\chapterline

\section{Notation and conventions}
\label{ch40:notation}

Throughout, $\Z = \set{\dots,-2,-1,0,1,2,\dots}$ is the integers and $\N = \set{0,1,2,\dots}$ the
non-negative integers. Lowercase letters $a,b,c,m,n,p,q$ denote integers; $p$ and $q$ are almost
always prime. We write:

\begin{itemize}
\item $a \mid b$ for ``$a$ divides $b$'': there is an integer $k$ with $b = ak$.
\item $\gcd(a,b)$ for the greatest common divisor, and $\lcm(a,b)$ for the least common multiple.
\item $a \equiv b \pmod{n}$ for ``$n \mid (a-b)$'', read ``$a$ is congruent to $b$ modulo $n$''.
\item $\Zn{n}$ for the ring of residues $\set{0,1,\dots,n-1}$ under addition and multiplication
      mod $n$.
\item $\Zns{n}$ for the group of \emph{units} mod $n$ --- those residues $a$ with $\gcd(a,n)=1$ ---
      under multiplication.
\item $\varphi(n)$ for Euler's totient, the size of $\Zns{n}$.
\item $\ord_n(a)$ for the multiplicative order of $a$ modulo $n$.
\item $a^{-1} \bmod n$ for the multiplicative inverse of $a$ modulo $n$, when it exists.
\end{itemize}

A residue and its representative are conflated freely: ``$3 \in \Zn{7}$'' means the class of all
integers congruent to $3$ mod $7$, but we name it by its smallest non-negative member. When we say
``reduce mod $n$'' we mean replace an integer by the unique representative in $\set{0,\dots,n-1}$.

\begin{intuitionbox}
Modular arithmetic is clock arithmetic. On a 12-hour clock, $7 + 8 = 15$ becomes $3$ because you
wrap around at $12$. Everything in this appendix is the consequence of one idea: once you only care
about remainders after dividing by $n$, the integers collapse into a finite world of exactly $n$
values, and that finite world is rich enough to hide secrets in.
\end{intuitionbox}

\section{Divisibility and the division algorithm}
\label{ch40:division}

\begin{theorem}[Division algorithm]
\label{ch40:thm-division}
For any integer $a$ and any positive integer $n$, there exist unique integers $q$ (quotient) and
$r$ (remainder) with
\[
  a = qn + r, \qquad 0 \le r < n.
\]
\end{theorem}

The remainder $r$ is what we call $a \bmod n$. Existence is just ``keep subtracting $n$ until you
land in $[0,n)$''; uniqueness follows because two remainders differing by less than $n$ whose
difference is divisible by $n$ must be equal. Worked: $a = 47$, $n = 9$ gives $47 = 5\cdot 9 + 2$, so
$47 \bmod 9 = 2$. For negatives we still demand $r \ge 0$: $-47 = (-6)\cdot 9 + 7$, so
$-47 \bmod 9 = 7$, \emph{not} $-2$. This sign convention matters in code; see the trap below.

\begin{trapbox}
In C, Java, and JavaScript the \texttt{\%} operator follows the sign of the \emph{dividend}, so
\texttt{-47 \% 9} returns $-2$, not the mathematician's $7$. Cryptographic code that indexes tables
or computes Legendre symbols with a raw \texttt{\%} on a possibly-negative value will silently read
the wrong slot. The fix is \texttt{((a \% n) + n) \% n}, or use a big-integer library whose
\texttt{mod} already normalizes to $[0,n)$. Python's \texttt{\%} already does the right thing.
\end{trapbox}

\section{Greatest common divisor and the Euclidean algorithm}
\label{ch40:gcd}

The greatest common divisor $\gcd(a,b)$ is the largest integer dividing both $a$ and $b$ (with
$\gcd(0,0)$ left undefined and $\gcd(a,0)=\abs{a}$). Two numbers with $\gcd = 1$ are
\emph{coprime}. The Euclidean algorithm computes the gcd using one fact:
\[
  \gcd(a,b) = \gcd(b,\ a \bmod b).
\]
Any common divisor of $a$ and $b$ divides $a - qb = a\bmod b$, and vice versa, so the set of common
divisors never changes as we replace $(a,b)$ by $(b, a\bmod b)$; the second coordinate shrinks
every step until it hits $0$, and then the first coordinate is the gcd.

\begin{example}[Euclid on $1071$ and $462$]
\label{ch40:ex-euclid}
\[
\begin{aligned}
\gcd(1071,462) &= \gcd(462,\ 1071 \bmod 462) = \gcd(462, 147)\\
               &= \gcd(147,\ 462 \bmod 147) = \gcd(147, 21)\\
               &= \gcd(21,\ 147 \bmod 21)  = \gcd(21, 0) = 21.
\end{aligned}
\]
\end{example}

\begin{connectionbox}
The Euclidean algorithm is very fast: the number of steps is $O(\log \min(a,b))$, and the
worst case is consecutive Fibonacci numbers (Lam\'e's theorem). The companion algorithms book
analyzes this; the upshot for cryptography is that gcd, modular inverse, and therefore RSA key
generation are all cheap even on 4096-bit numbers. Contrast that with factoring those same numbers,
which nobody knows how to do quickly --- the whole field lives in that asymmetry.
\end{connectionbox}

\subsection{B\'ezout and the extended Euclidean algorithm}
\label{ch40:bezout}

\begin{theorem}[B\'ezout's identity]
\label{ch40:thm-bezout}
For integers $a,b$ not both zero, there exist integers $s,t$ (the \emph{B\'ezout coefficients})
with
\[
  as + bt = \gcd(a,b).
\]
\end{theorem}

The extended Euclidean algorithm finds $s$ and $t$ alongside the gcd by carrying the combination
through every step. This is the single most useful subroutine in public-key crypto, because the
case $\gcd(a,n)=1$ gives $as + nt = 1$, i.e.\ $as \equiv 1 \pmod n$, so $s = a^{-1} \bmod n$.

\begin{example}[Extended Euclid: inverse of $7$ mod $26$]
\label{ch40:ex-extended}
We track triples $(r,s,t)$ with $r = 7s + 26t$. Start from $(26,0,1)$ and $(7,1,0)$; each row
subtracts $q$ times the next.
\[
\begin{array}{c|c|c|c}
\text{remainder } r & \text{quotient } q & s & t\\\hline
26 &   & 0 & 1\\
7  &   & 1 & 0\\
5  & 3 & -3 & 1\\
2  & 1 & 4 & -1\\
1  & 2 & -11 & 3\\
0  & 2 &  &
\end{array}
\]
The last nonzero remainder is $1 = \gcd(7,26)$, with $s=-11,\ t=3$: check $7\cdot(-11) + 26\cdot 3
= -77 + 78 = 1$. So $7^{-1} \equiv -11 \equiv 15 \pmod{26}$, and indeed $7\cdot 15 = 105 = 4\cdot 26
+ 1$.
\end{example}

\begin{implementbox}
A compact iterative extended-gcd in Python. It returns $(g,s,t)$ with $as+bt=g$.
\begin{lstlisting}[language=Python]
def ext_gcd(a, b):
    old_r, r = a, b
    old_s, s = 1, 0
    old_t, t = 0, 1
    while r != 0:
        q = old_r // r
        old_r, r = r, old_r - q * r
        old_s, s = s, old_s - q * s
        old_t, t = t, old_t - q * t
    return old_r, old_s, old_t          # g, s, t  with a*s + b*t = g

def inverse(a, n):
    g, s, _ = ext_gcd(a % n, n)
    if g != 1:
        raise ValueError("no inverse: gcd != 1")
    return s % n                        # normalize into [0, n)

# inverse(7, 26) -> 15
\end{lstlisting}
LEARNING ONLY. Real libraries use constant-time, side-channel-hardened modular inverses (e.g.\
Fermat-based or a binary GCD with branchless updates) because the loop count above leaks
information about the secret operand through timing. Never ship this for key material.
\end{implementbox}

\section{Modular arithmetic: the ring $\Zn{n}$}
\label{ch40:modring}

Congruence respects addition, subtraction, and multiplication: if $a\equiv a'$ and $b\equiv b'$
$\pmod n$, then $a\pm b \equiv a'\pm b'$ and $ab\equiv a'b'\pmod n$. So we can reduce ``early and
often'' --- replace any intermediate result by its remainder without changing the final answer.
This is why modular exponentiation of huge numbers never overflows: you reduce after each multiply.

\begin{walkbox}
Why can't we always divide? In $\Zn{n}$, division by $a$ means multiplying by $a^{-1}$, and that
inverse exists \emph{iff} $\gcd(a,n)=1$. In $\Zn{12}$ the element $3$ has no inverse: $3\cdot k$ runs
through $\set{0,3,6,9}$ and never hits $1$. Trace it: $3\cdot 1=3$, $3\cdot 2=6$, $3\cdot 3=9$,
$3\cdot 4=12\equiv 0$, then it repeats. The multiples of a non-coprime element are trapped on a
sublattice. Only the units --- the residues coprime to $n$ --- are invertible.
\end{walkbox}

\begin{definition}[Units]
\label{ch40:def-units}
$\Zns{n} = \set{a \in \Zn{n} : \gcd(a,n)=1}$ is the group of units under multiplication mod $n$. It
is closed, associative, has identity $1$, and every element has an inverse (by B\'ezout). When $n=p$
is prime, every nonzero residue is a unit, so $\Zns{p}=\set{1,2,\dots,p-1}$ and $\Zn{p}$ is a field.
\end{definition}

\subsection{Fast modular exponentiation}
\label{ch40:fastexp}

Computing $a^e \bmod n$ by multiplying $a$ by itself $e$ times is hopeless when $e$ has hundreds of
digits. Square-and-multiply does it in $O(\log e)$ multiplications by reading the exponent's binary
digits. To compute $a^{13}\bmod n$, note $13 = 1101_2$, so
$a^{13} = a^{8}\cdot a^{4}\cdot a^{1}$: square repeatedly to get $a,a^2,a^4,a^8$ and multiply in the
ones that the bits select.

\begin{example}[$3^{13}\bmod 7$]
\label{ch40:ex-modexp}
Repeated squaring mod $7$: $3^1=3$, $3^2=9\equiv 2$, $3^4\equiv 2^2=4$, $3^8\equiv 4^2=16\equiv 2$.
Bits of $13=1101_2$ pick $a^8,a^4,a^1$: $3^{13}\equiv 2\cdot 4\cdot 3 = 24 \equiv 3 \pmod 7$. (Sanity
check via Fermat below: $3^6\equiv 1$, so $3^{13}=3^{12}\cdot 3\equiv 1\cdot 3 = 3$.)
\end{example}

\begin{implementbox}
\begin{lstlisting}[language=Python]
def mod_exp(a, e, n):
    result = 1
    a = a % n
    while e > 0:
        if e & 1:                 # current low bit is set
            result = (result * a) % n
        a = (a * a) % n           # square the base
        e >>= 1                   # shift to next bit
    return result
# mod_exp(3, 13, 7) -> 3
\end{lstlisting}
LEARNING ONLY. The branch \texttt{if e \& 1} makes the running time depend on the secret exponent's
bits --- a classic timing/power side channel that has broken real RSA and ECC implementations.
Production code uses Montgomery ladders or constant-time windowed exponentiation. Python's built-in
\texttt{pow(a, e, n)} is the right call for learning, and is itself not constant-time.
\end{implementbox}

\section{Fermat's little theorem and Euler's theorem}
\label{ch40:fermateuler}

\begin{theorem}[Fermat's little theorem]
\label{ch40:thm-fermat}
If $p$ is prime and $p \nmid a$, then
\[
  a^{p-1} \equiv 1 \pmod p.
\]
Equivalently $a^p \equiv a \pmod p$ for \emph{every} integer $a$.
\end{theorem}

\begin{whybox}
Here is the one-line proof worth remembering. Multiply the nonzero residues
$\set{1,2,\dots,p-1}$ each by $a$. Since $a$ is a unit, the products $\set{a\cdot 1,\dots,a\cdot
(p-1)}$ are the same set, just shuffled. Multiply both sets together:
\[
  (p-1)! \equiv a^{p-1}\,(p-1)! \pmod p.
\]
Now $(p-1)!$ is coprime to $p$, so cancel it: $a^{p-1}\equiv 1$. The whole theorem is ``relabeling
a set does not change its product.''
\end{whybox}

\begin{theorem}[Euler's theorem]
\label{ch40:thm-euler}
If $\gcd(a,n)=1$ then
\[
  a^{\varphi(n)} \equiv 1 \pmod n,
\]
where $\varphi(n) = \abs{\Zns{n}}$ is Euler's totient.
\end{theorem}

The proof is identical: multiplication by a unit $a$ permutes $\Zns{n}$, so the product of all units
equals $a^{\varphi(n)}$ times itself, and cancellation gives the result. Fermat is the special case
$n=p$, where $\varphi(p)=p-1$.

\subsection{Computing the totient}
\label{ch40:totient}

$\varphi$ is multiplicative on coprime arguments and easy on prime powers:
\[
  \varphi(p) = p-1,\qquad
  \varphi(p^k) = p^k - p^{k-1} = p^{k}\Bigl(1-\tfrac1p\Bigr),\qquad
  \varphi(mn) = \varphi(m)\varphi(n)\ \text{ if } \gcd(m,n)=1.
\]
So from the prime factorization $n = \prod_i p_i^{k_i}$,
\[
  \varphi(n) = n\prod_{i}\Bigl(1 - \tfrac{1}{p_i}\Bigr).
\]

\begin{example}[$\varphi$ of a few numbers]
\label{ch40:ex-totient}
$\varphi(10)=\varphi(2)\varphi(5)=1\cdot 4 = 4$ (the units are $1,3,7,9$). $\varphi(36)=\varphi(4)
\varphi(9)=2\cdot 6 = 12$. For an RSA-style modulus $n=pq$ with distinct primes, $\varphi(n) =
(p-1)(q-1)$: e.g.\ $p=61,q=53$ give $n=3233$ and $\varphi(n)=60\cdot 52 = 3120$.
\end{example}

\begin{applicationbox}
Euler's theorem is the engine of RSA correctness. With $n=pq$ and exponents $e,d$ chosen so that
$ed\equiv 1\pmod{\varphi(n)}$, decryption recovers the message because $m^{ed} = m^{1 + k\varphi(n)}
= m\cdot(m^{\varphi(n)})^k \equiv m\pmod n$. The private key $d$ is literally $e^{-1}\bmod
\varphi(n)$, computed with the extended Euclidean algorithm of \S\ref{ch40:bezout}. Knowing
$\varphi(n)$ is equivalent to factoring $n$ --- which is why $p,q$ must stay secret.
\end{applicationbox}

\section{The Chinese Remainder Theorem}
\label{ch40:crt}

\begin{theorem}[CRT]
\label{ch40:thm-crt}
Let $n_1,\dots,n_k$ be pairwise coprime and $N = \prod_i n_i$. Then the system
\[
  x \equiv a_1 \pmod{n_1},\quad \dots,\quad x \equiv a_k \pmod{n_k}
\]
has a unique solution modulo $N$. Equivalently, the map $x \mapsto (x\bmod n_1,\dots,x\bmod n_k)$ is
a ring isomorphism $\Zn{N} \cong \Zn{n_1}\times\cdots\times\Zn{n_k}$.
\end{theorem}

The isomorphism view is the one to internalize: arithmetic mod $N$ \emph{is} component-wise
arithmetic on the small moduli. Add, multiply, exponentiate each coordinate independently, then glue
back. The gluing is explicit. For two moduli, write $N=n_1 n_2$ and find $m_1 = n_2^{-1}\bmod n_1$,
$m_2 = n_1^{-1}\bmod n_2$; then
\[
  x = a_1\, n_2\, m_1 + a_2\, n_1\, m_2 \pmod N.
\]

\begin{example}[Solve $x\equiv 2\ (5)$, $x\equiv 3\ (7)$]
\label{ch40:ex-crt}
Here $n_1=5,n_2=7,N=35$. Compute $m_1 = 7^{-1}\bmod 5$: $7\equiv 2$, and $2\cdot 3=6\equiv 1$, so
$m_1=3$. Compute $m_2 = 5^{-1}\bmod 7$: $5\cdot 3=15\equiv 1$, so $m_2=3$. Then
\[
  x = 2\cdot 7\cdot 3 + 3\cdot 5\cdot 3 = 42 + 45 = 87 \equiv 87 - 2\cdot 35 = 17 \pmod{35}.
\]
Check: $17\bmod 5 = 2$ and $17\bmod 7 = 3$. Correct.
\end{example}

The isomorphism $\Zn{35}\cong\Zn{5}\times\Zn{7}$ says each $x\in\set{0,\dots,34}$ sits at a unique
cell of a $5\times 7$ grid indexed by its two remainders, and every cell is hit exactly once.
Figure~\ref{ch40:fig-crt} plots a few values; the pair $(2,3)$ pins down $x=17$.

\begin{figure}[H]
\centering
\begin{tikzpicture}[x=0.95cm,y=0.95cm]
  % axes: columns = residue mod 7 (0..6), rows = residue mod 5 (0..4)
  \draw[gray!55] (0,0) grid (7,5);
  \node at (3.5,-0.8) {residue mod $7$ $\longrightarrow$};
  \node[rotate=90] at (-0.8,2.5) {residue mod $5$ $\longrightarrow$};
  \foreach \c in {0,...,6} { \node[below] at (\c+0.5,0) {\footnotesize \c}; }
  \foreach \r in {0,...,4} { \node[left]  at (0,\r+0.5) {\footnotesize \r}; }
  % place each x at (x mod 7 + .5, x mod 5 + .5); highlight 17
  \foreach \x in {0,...,34} {
    \pgfmathtruncatemacro{\cc}{mod(\x,7)}
    \pgfmathtruncatemacro{\rr}{mod(\x,5)}
    \ifnum\x=17
      \node[fill=trapred,text=white,circle,inner sep=1pt,font=\scriptsize\bfseries]
        at (\cc+0.5,\rr+0.5) {\x};
    \else
      \node[font=\tiny,gray!75] at (\cc+0.5,\rr+0.5) {\x};
    \fi
  }
\end{tikzpicture}
\caption{The CRT bijection $\Zn{35}\cong\Zn{5}\times\Zn{7}$. Each residue mod $35$ lands in one cell,
keyed by its remainders mod $5$ (row) and mod $7$ (column); the system $x\equiv 2\ (5),\ x\equiv 3\
(7)$ selects the highlighted cell $x=17$.}
\label{ch40:fig-crt}
\end{figure}

\begin{applicationbox}
RSA decryption and signing run about four times faster with CRT. Instead of one exponentiation mod
$n$ (a $2k$-bit number), you do two exponentiations mod $p$ and mod $q$ (each $k$-bit), then glue.
The cost of an exponentiation grows roughly with the cube of the bit length, so halving the modulus
size cuts each piece to $1/8$, and two of them total $1/4$. OpenSSL, BoringSSL, and every serious
RSA implementation store the CRT parameters $d_p = d\bmod(p-1)$, $d_q = d\bmod(q-1)$, and
$q^{-1}\bmod p$ in the private key for exactly this reason.
\end{applicationbox}

\begin{attackbox}
CRT-RSA has a notorious failure mode. If a hardware fault (a glitch, a cosmic ray, a deliberate
laser) corrupts the computation modulo $p$ but not modulo $q$, the faulty signature $\hat s$
satisfies $\hat s^e \equiv m \pmod q$ but $\hat s^e \not\equiv m\pmod p$. Then
\[
  \gcd(\hat s^e - m,\ n) = q,
\]
handing the attacker a prime factor and thus the entire private key from \emph{one} faulty
signature. This is the Boneh-DeMillo-Lipton fault attack. Defenses verify the signature (recompute
$\hat s^e \bmod n$ and compare to $m$) before releasing it.
\end{attackbox}

\section{Orders, cyclic groups, and primitive roots}
\label{ch40:orders}

\begin{definition}[Multiplicative order]
\label{ch40:def-order}
For $a\in\Zns{n}$, the \emph{order} $\ord_n(a)$ is the least positive integer $d$ with $a^d\equiv 1
\pmod n$.
\end{definition}

By Euler's theorem $a^{\varphi(n)}\equiv 1$, and a short argument shows $\ord_n(a)\mid\varphi(n)$:
divide $\varphi(n)$ by $d$ with remainder $r$, and $a^{r}\equiv 1$ would follow, contradicting
minimality of $d$ unless $r=0$. So orders are always divisors of the group size --- a special case of
Lagrange's theorem.

\begin{example}[Orders in $\Zns{7}$]
\label{ch40:ex-order}
Here $\varphi(7)=6$, so every order divides $6$. Powers of $3$: $3,2,6,4,5,1$ --- order $6$. Powers
of $2$: $2,4,1$ --- order $3$. Powers of $6$: $6,1$ --- order $2$. The possible orders $1,2,3,6$ all
appear, and $3$ runs through \emph{all} six units, so $3$ is a generator.
\end{example}

\begin{definition}[Primitive root]
\label{ch40:def-primroot}
A \emph{primitive root} modulo $n$ is an element $g\in\Zns{n}$ of order $\varphi(n)$; its powers
generate the whole group, making $\Zns{n}$ \emph{cyclic} with generator $g$.
\end{definition}

\begin{theorem}[Existence of primitive roots]
\label{ch40:thm-primroot}
$\Zns{n}$ is cyclic --- a primitive root exists --- exactly when $n\in\set{1,2,4,p^k,2p^k}$ for an
odd prime $p$. In particular every prime modulus $p$ has a primitive root, and $\Zns{p}$ is cyclic
of order $p-1$.
\end{theorem}

That last sentence is what Diffie-Hellman and ElGamal stand on: a prime $p$ gives a cyclic group in
which exponentiation is a permutation but its inverse (the discrete logarithm) is believed hard.

\begin{whybox}
A cyclic group of order $m$ has exactly $\varphi(m)$ generators, and for each divisor $d\mid m$
exactly $\varphi(d)$ elements of order $d$. Why $\varphi(m)$ generators? If $g$ generates, then
$g^k$ generates iff $\gcd(k,m)=1$ --- otherwise $g^k$ cycles through a proper subgroup. In $\Zns{7}$
($m=6$) there are $\varphi(6)=2$ primitive roots, namely $3$ and $5$, matching the count from the
example above.
\end{whybox}

\begin{connectionbox}
To \emph{test} whether $g$ is a primitive root mod $p$ you do not check all $p-2$ powers. Factor
$p-1 = \prod q_i^{e_i}$ and verify $g^{(p-1)/q_i}\not\equiv 1\pmod p$ for each prime $q_i$. If none of
those drops to $1$, then $\ord_p(g)$ cannot be a proper divisor of $p-1$, so it equals $p-1$. This is
why ``safe primes'' $p = 2q+1$ (with $q$ prime) are popular: $p-1 = 2q$ has only two prime factors,
so generator testing is trivial and the only subgroups have orders $1,2,q,2q$.
\end{connectionbox}

\section{Primality testing}
\label{ch40:primality}

Generating RSA or Diffie-Hellman parameters means producing large primes --- often thousands of
bits. Trial division is hopeless past a few digits. Probabilistic tests built on Fermat's theorem
are the workhorses.

\subsection{The Fermat test and its blind spot}
\label{ch40:fermat-test}

If $n$ is prime, Fermat says $a^{n-1}\equiv 1\pmod n$ for all $a$ coprime to $n$. So pick a random
base $a$; if $a^{n-1}\not\equiv 1$, then $n$ is \emph{definitely composite} ($a$ is a ``witness'').
If it equals $1$, $n$ is probably prime. The blind spot:

\begin{attackbox}
\emph{Carmichael numbers} are composite $n$ for which $a^{n-1}\equiv 1\pmod n$ for \emph{every} $a$
coprime to $n$ --- the Fermat test never finds a witness no matter how many bases you try. The
smallest is $561 = 3\cdot 11\cdot 17$. There are infinitely many. A naive Fermat-only primality
check will cheerfully accept $561$ as prime, and an RSA modulus built on such a ``prime'' factors
instantly. This is why production code uses Miller-Rabin, which has no Carmichael blind spot.
\end{attackbox}

\subsection{Miller-Rabin}
\label{ch40:millerrabin}

Miller-Rabin strengthens the Fermat test by looking inside the squaring chain. Write $n-1 =
2^{s}\,d$ with $d$ odd. For a base $a$, compute $x = a^d \bmod n$ and then square it $s$ times. If
$n$ is prime, the sequence $a^d, a^{2d}, \dots, a^{2^s d}=a^{n-1}$ must end at $1$, and the first time
it reaches $1$ it must do so \emph{from} $-1$ (because in a field the only square roots of $1$ are
$\pm 1$). So $n$ \emph{passes} for base $a$ if either
\[
  a^d \equiv 1 \pmod n,\qquad\text{or}\qquad a^{2^r d}\equiv -1\pmod n\ \text{ for some }0\le r<s.
\]
Otherwise $a$ is a witness and $n$ is composite. The key theorem: for an odd composite $n$, at least
$3/4$ of the bases in $\set{1,\dots,n-1}$ are witnesses. So $k$ independent random bases that all
pass leave a failure probability of at most $4^{-k}$.

\begin{example}[Miller-Rabin catches $561$]
\label{ch40:ex-mr}
$n=561$, so $n-1 = 560 = 2^4\cdot 35$, giving $s=4,\ d=35$. Take base $a=2$. Compute $2^{35}\bmod
561 = 263$ (not $\pm 1$). Square: $263^2 = 69169 \equiv 166\pmod{561}$. Square: $166^2 = 27556\equiv
67$. Square: $67^2 = 4489\equiv 1\pmod{561}$. We reached $1$, but the value just before it was $67$,
not $560\equiv -1$. A nontrivial square root of $1$ appeared --- impossible in a field --- so $2$ is a
witness and $561$ is composite. The Fermat test missed this because $2^{560}\equiv 1$; Miller-Rabin
caught the illegal square root on the way up.
\end{example}

\begin{implementbox}
\begin{lstlisting}[language=Python]
import random

def miller_rabin(n, rounds=40):
    if n < 2:            return False
    for p in (2,3,5,7,11,13,17,19,23,29,31,37):
        if n % p == 0:   return n == p
    # write n-1 = 2^s * d  with d odd
    d, s = n - 1, 0
    while d % 2 == 0:
        d //= 2; s += 1
    for _ in range(rounds):
        a = random.randrange(2, n - 1)
        x = pow(a, d, n)
        if x == 1 or x == n - 1:
            continue
        for _ in range(s - 1):
            x = (x * x) % n
            if x == n - 1:
                break
        else:
            return False           # a is a witness: composite
    return True                    # probably prime
\end{lstlisting}
LEARNING ONLY. This uses non-cryptographic \texttt{random} and is not constant-time. For real key
generation, draw bases from a CSPRNG and follow standards (FIPS 186-5 specifies round counts by key
size); libraries also combine Miller-Rabin with a Lucas test (the Baillie-PSW combination) so that no
known composite passes.
\end{implementbox}

\begin{historybox}
A deterministic polynomial-time primality test (AKS, 2002, by Agrawal, Kayal, and Saxena) settled the
theoretical question --- PRIMES is in P --- but it is far slower in practice than Miller-Rabin. Under
the Generalized Riemann Hypothesis, testing only bases up to $2(\ln n)^2$ makes Miller-Rabin
deterministic; cryptographers instead just run enough random rounds that the error is astronomically
below hardware-fault rates and move on.
\end{historybox}

\section{Quadratic residues}
\label{ch40:qr}

\begin{definition}[Quadratic residue]
\label{ch40:def-qr}
For odd prime $p$ and $a\not\equiv 0$, $a$ is a \emph{quadratic residue} (QR) mod $p$ if $x^2\equiv a
\pmod p$ has a solution; otherwise a \emph{quadratic non-residue} (QNR).
\end{definition}

Exactly half of the nonzero residues are QRs: the squaring map $x\mapsto x^2$ is two-to-one (since
$x$ and $-x$ share a square), so its image has size $(p-1)/2$. Modulo $7$: the squares of
$1,\dots,6$ are $1,4,2,2,4,1$, so the QRs are $\set{1,2,4}$ and the QNRs are $\set{3,5,6}$.

\begin{definition}[Legendre symbol]
\label{ch40:def-legendre}
\[
  \left(\frac{a}{p}\right) =
  \begin{cases}
    \ \ 1 & a \text{ is a QR mod } p,\\
    -1 & a \text{ is a QNR mod } p,\\
    \ \ 0 & p\mid a.
  \end{cases}
\]
\end{definition}

\begin{theorem}[Euler's criterion]
\label{ch40:thm-eulercrit}
For odd prime $p$ and $p\nmid a$,
\[
  \left(\frac{a}{p}\right) \equiv a^{(p-1)/2} \pmod p.
\]
\end{theorem}

Why: $a^{(p-1)/2}$ squares to $a^{p-1}\equiv 1$, so it is a square root of $1$, hence $\pm 1$. It is
$+1$ precisely when $a$ is a square, because if $a=g^{2j}$ (a generator power with even exponent) then
$a^{(p-1)/2}=g^{j(p-1)}\equiv 1$, while odd exponents give $-1$. This is a fast QR test: one modular
exponentiation. Check mod $7$: $3^{3}=27\equiv 6\equiv -1$, so $3$ is a QNR (matches the list);
$2^{3}=8\equiv 1$, so $2$ is a QR.

\begin{example}[Computing $\left(\tfrac{2}{7}\right)$ and a square root]
\label{ch40:ex-sqrt}
$2$ is a QR mod $7$. When $p\equiv 3\pmod 4$ --- and $7$ is --- a square root is just $a^{(p+1)/4}$:
here $2^{(7+1)/4}=2^2=4$, and indeed $4^2=16\equiv 2\pmod 7$. The other root is $-4\equiv 3$; check
$3^2 = 9\equiv 2$. For $p\equiv 1\pmod 4$ you need the Tonelli-Shanks algorithm instead.
\end{example}

\begin{applicationbox}
Quadratic residues are everywhere. Elliptic-curve \emph{point compression} stores only the
$x$-coordinate plus one sign bit: the receiver recovers $y$ by taking a modular square root of
$x^3+ax+b$, which exists only when that value is a QR --- the sign bit picks which root. The
\emph{Goldwasser-Micali} cryptosystem encrypts a $0$ bit as a random QR and a $1$ bit as a
QNR modulo $n=pq$, and its security rests on the \emph{quadratic residuosity problem}: deciding
QR-ness modulo a composite without the factorization is believed hard. \emph{Rabin encryption}
squares the message mod $n$, and decrypting means taking square roots mod $p$ and $q$ and CRT-gluing
the four roots. And the Jacobi symbol (the multiplicative extension of Legendre to composite moduli)
lets you compute these symbols fast \emph{without} factoring, via quadratic reciprocity.
\end{applicationbox}

\begin{trapbox}
The Jacobi symbol $\left(\tfrac{a}{n}\right)$ for composite $n$ can equal $+1$ even when $a$ is a
\emph{non}-residue mod $n$ (it only tracks the product of Legendre symbols, so two $-1$'s cancel). So
a Jacobi value of $+1$ does \emph{not} prove $a$ is a square modulo a composite. Goldwasser-Micali
exploits exactly this gap: it hides bits inside the pseudosquares (Jacobi $+1$ but not a true QR),
which are indistinguishable from real squares without the factorization.
\end{trapbox}

\section{Quick reference}
\label{ch40:tables}

\begin{focusbox}
\textbf{Core identities.}
\[
\begin{aligned}
&\gcd(a,b) = as+bt && \text{(B\'ezout; extended Euclid finds }s,t)\\
&a^{-1}\bmod n \text{ exists} \iff \gcd(a,n)=1 && (s\bmod n \text{ from above})\\
&a^{p-1}\equiv 1 \pmod p && \text{(Fermat, }p\nmid a)\\
&a^{\varphi(n)}\equiv 1 \pmod n && \text{(Euler, }\gcd(a,n)=1)\\
&\varphi(n) = n\textstyle\prod_{p\mid n}(1-\tfrac1p) && \text{(totient from factorization)}\\
&\ord_n(a) \mid \varphi(n) && \text{(orders divide the group size)}\\
&\Zn{N}\cong \textstyle\prod_i \Zn{n_i} && \text{(CRT, }n_i\text{ pairwise coprime)}\\
&\left(\tfrac{a}{p}\right)\equiv a^{(p-1)/2}\pmod p && \text{(Euler's criterion)}
\end{aligned}
\]
\end{focusbox}

\begin{center}
\renewcommand{\arraystretch}{1.25}
\begin{tabular}{@{}lll@{}}
\toprule
\textbf{Task} & \textbf{Tool} & \textbf{Cost}\\
\midrule
$\gcd(a,b)$                  & Euclidean algorithm          & $O(\log\min(a,b))$ steps\\
$a^{-1}\bmod n$              & extended Euclid              & $O(\log n)$ steps\\
$a^e \bmod n$               & square-and-multiply          & $O(\log e)$ multiplications\\
system of $x\bmod n_i$       & CRT reconstruction           & $O(k)$ inverses\\
``is $n$ prime?''            & Miller-Rabin ($k$ rounds)    & error $\le 4^{-k}$\\
``is $a$ a square mod $p$?'' & Euler's criterion            & one exponentiation\\
$\sqrt{a}\bmod p$, $p\equiv 3\ (4)$ & $a^{(p+1)/4}$         & one exponentiation\\
\bottomrule
\end{tabular}
\end{center}

\section{Worked end-to-end: a complete RSA toy with these tools}
\label{ch40:rsa-toy}

To see every piece fire at once, run RSA with $p=61,\ q=53$ (the classic textbook numbers).

\begin{enumerate}
\item \textbf{Modulus and totient.} $n = 61\cdot 53 = 3233$, and $\varphi(n) = 60\cdot 52 = 3120$
      (totient of a product, \S\ref{ch40:totient}).
\item \textbf{Public exponent.} Choose $e = 17$; check $\gcd(17,3120)=1$ with Euclid.
\item \textbf{Private exponent.} $d = 17^{-1}\bmod 3120$ via extended Euclid (\S\ref{ch40:bezout}).
      Running it gives $d = 2753$; verify $17\cdot 2753 = 46801 = 15\cdot 3120 + 1\equiv 1$.
\item \textbf{Encrypt} $m = 65$: $c = 65^{17}\bmod 3233$ by square-and-multiply $= 2790$.
\item \textbf{Decrypt} with CRT for speed (\S\ref{ch40:crt}). Compute
      $d_p = d\bmod (p-1) = 2753\bmod 60 = 53$ and $d_q = d\bmod(q-1) = 2753\bmod 52 = 49$. Then
      \[
        m_p = c^{d_p}\bmod 61 = 2790^{53}\bmod 61 = 4,\qquad
        m_q = c^{d_q}\bmod 53 = 2790^{49}\bmod 53 = 12.
      \]
      CRT-glue $x\equiv 4\ (61),\ x\equiv 12\ (53)$: with $61^{-1}\bmod 53 = 20$ and $53^{-1}\bmod
      61 = 38$, reconstruct $m = 4\cdot 53\cdot 38 + 12\cdot 61\cdot 20 = 8056 + 14640 = 22696
      \equiv 65 \pmod{3233}$. The Euler-theorem identity $m^{ed}\equiv m$ guarantees this.
\end{enumerate}

Every line used a tool from this appendix: factorization-based totient, extended Euclid for the
inverse, fast exponentiation, CRT for the speedup. That is the whole appendix in miniature, and it is
exactly what your TLS library does --- just with $2048$-bit primes instead of two-digit ones.

\section{Exercises}
\label{ch40:exercises}

\begin{exercise}
\label{ch40:ex-gcd}
Use the Euclidean algorithm to compute $\gcd(2024, 748)$, showing each remainder. Then run extended
Euclid to find integers $s,t$ with $2024s + 748t = \gcd$.
\end{exercise}

\begin{exercise}
\label{ch40:ex-inverse}
Compute $11^{-1}\bmod 26$ two ways: (a) by extended Euclid, and (b) by Euler's theorem, using
$11^{\varphi(26)-1}\bmod 26$. Confirm both give the same residue, and state $\varphi(26)$.
\end{exercise}

\begin{exercise}
\label{ch40:ex-modexp2}
Compute $7^{222}\bmod 11$. First reduce the exponent using Fermat's little theorem, then finish by
hand. (Hint: what is $7^{10}\bmod 11$?)
\end{exercise}

\begin{exercise}
\label{ch40:ex-crt2}
Solve the system $x\equiv 1\pmod 3$, $x\equiv 2\pmod 4$, $x\equiv 3\pmod 5$ for the unique
$x\in\set{0,\dots,59}$. Verify your answer against all three congruences.
\end{exercise}

\begin{exercise}
\label{ch40:ex-order2}
List the order of every element of $\Zns{11}$. How many primitive roots does $11$ have, and does
this match the $\varphi(\varphi(11))$ prediction? Exhibit one primitive root.
\end{exercise}

\begin{exercise}
\label{ch40:ex-carmichael}
Verify that $561 = 3\cdot 11\cdot 17$ is a Carmichael number by checking the Korselt criterion: $n$
is squarefree and $(p-1)\mid(n-1)$ for each prime factor $p$. Then find a single Miller-Rabin witness
for $561$ other than the base $2$ used in Example~\ref{ch40:ex-mr}.
\end{exercise}

\begin{exercise}
\label{ch40:ex-qr}
Determine all quadratic residues modulo $13$ by squaring $1,\dots,12$. Then use Euler's criterion to
confirm whether $5$ and $6$ are residues, and for each residue among them give a square root.
\end{exercise}

\begin{exercise}
\label{ch40:ex-fault}
(Break it.) An RSA signer uses CRT and a fault corrupts the mod-$p$ branch, producing a signature
$\hat s$ with $\hat s^e\equiv m\pmod q$ but $\hat s^e\not\equiv m\pmod p$. Explain why
$\gcd(\hat s^e - m,\ n)$ recovers the factor $q$, and why verifying the signature before output
prevents the leak. Reference Theorem~\ref{ch40:thm-crt} in your reasoning.
\end{exercise}

\begin{exercise}
\label{ch40:ex-totient-prove}
(Prove it.) Show that if $\gcd(m,n)=1$ then $\varphi(mn)=\varphi(m)\varphi(n)$, using the CRT
isomorphism $\Zn{mn}\cong\Zn{m}\times\Zn{n}$ and the fact that an element is a unit iff it is a unit
in each coordinate.
\end{exercise}

\begin{exercise}
\label{ch40:ex-primroot-test}
Let $p = 23$, so $p-1 = 22 = 2\cdot 11$. Use the generator test of \S\ref{ch40:orders} (check
$g^{(p-1)/2}$ and $g^{(p-1)/11}$) to decide whether $g=5$ is a primitive root mod $23$, without
listing all $22$ powers.
\end{exercise}
